\begin{frame}{TCP Congestion Window Overview}

\begin{center}
\textbf{$cwnd$ is the sender's estimate of how much traffic the network can currently sustain.}
\end{center}

\[
\boxed{Sending\ rate \approx \frac{cwnd}{RTT}}
\]

\begin{columns}[T,onlytextwidth]
\column{0.47\textwidth}
\textbf{Core role}
\begin{itemize}
    \item Maintained by the TCP sender.
    \item Controls the amount of outstanding data.
    \item Separate from the receiver window $rwnd$.
\end{itemize}

\column{0.49\textwidth}
\textbf{Feedback-driven adaptation}
\begin{itemize}
    \item Successful ACKs allow $cwnd$ to grow.
    \item Loss, timeout, or ECN causes $cwnd$ to decrease.
    \item The exact rule depends on the TCP algorithm.
\end{itemize}
\end{columns}

\vspace{0cm}
\begin{center}
\begin{tikzpicture}[node distance=0.55cm]
\node[draw=ResearchColor,rounded corners,fill=CardBg] (ack) {ACK feedback};
\node[draw=ResearchColor,rounded corners,fill=CardBg,right=of ack] (grow) {Increase $cwnd$};
\node[draw=WarningColor,rounded corners,fill=CardBg,right=of grow] (loss) {Congestion signal};
\node[draw=WarningColor,rounded corners,fill=CardBg,right=of loss] (reduce) {Reduce $cwnd$};
\draw[-{Latex[length=2mm]},draw=ResearchColor] (ack)--(grow);
\draw[-{Latex[length=2mm]},draw=WarningColor] (loss)--(reduce);
\end{tikzpicture}
\end{center}

\end{frame}
